% !TeX root = ../../AdvMacroNote_Sun.tex
\cleardoublepage
\pagestyle{symplain}
\thispagestyle{symplain}

\chapter*{第一版前言}
\addcontentsline{toc}{chapter}{第一版前言}
\vspace{\baselineskip}

% 公式编号仅用作引用，不代表重要程度

这是一份《高级宏观经济学 II》的笔记，由我于 2026 年春季学期在厦门大学修读该课程过程中所写。笔记内容由任课教师贺骋远的课程讲义（占大部分）、对讲义补充解释、个人思考、课外拓展构成。感谢你的翻阅，希望对你有些帮助。祝学习愉快～

\myskip
\mysec{笔记结构}
\begin{itemize}
    \item 前言；说明；凡例；目录
    \item \ref{part:03}
        \begin{itemize}
            \item \ref{ch:07}：主要围绕一阶、二阶线性差分系统展开，介绍如何求解方程和如何判定系统稳定性。
            \item \ref{ch:08}：经济模型通常是非线性的，难以求解，因此需要（对数）线性化方法；本章介绍多种线性化方法，是后续内容的基础工具。
        \end{itemize}
    \item \ref{part:04}
        \begin{itemize}
            \item \ref{ch:09}：根据实际经济周期（RBC）进行建模，探讨由消费、资本存量构成的经济系统，并判定其稳定性。
            \item \ref{ch:10}：继续 RBC 模型，致力于对模型进行数值求解，获得模型的数值解和动态路径。
        \end{itemize}
    \item 更多内容正在建设中……
    % \item 参考文献；后记；致谢
\end{itemize}

\newpage
\mysec{阅读建议}

\begin{itemize}
    \item 旁注是对正文的补充说明；脚注包含参考内容及勘误信息。
    \item 留意文档中的日期（如页眉和脚注），有助于识别更新和勘误。
    \item 版面使用按课程切换的主题强调色，并以正文、次要信息和分隔元素三个用途层级组织中性色。
    \item 附录位于每章末尾，包含拓展内容、过程细节。附录的序号中包含大写拉丁字母，如 2.C。
\end{itemize}

\myskip
\mysec{获取方式}

以下是获取该笔记的途径，进入仓库后请先阅读 \ttt{readme} 文件：
\begin{itemize}
    \item 厦大网盘：\url{https://box.xmu.edu.cn/share/037e63a9ff6d3dab1a680e26ef}
    \item GitHub：\url{https://github.com/econsun/XMU-AdvMacro-II}
\end{itemize}

\myskip
\mysec{触及方式}

目前这份笔记仍处于持续更新阶段，现阶段以跟进课堂进度为首要目标，缺少多轮校订，因此其中难免存在疏漏、错误或表述欠严谨之处，敬请指正，我会在你可接受范围内做好充分的致谢。请对笔记给予多些耐心，它会越来越好。

如果你想反馈问题、提出建议、交流学习，或参与笔记共建，欢迎通过以下任一方式触及我：
\begin{itemize}
    \item 小红书：\ttt{luckyEconSun}
    \item 邮箱：\ttt{econsunrq@outlook.com}
    \item 腾讯文档：\url{https://docs.qq.com/form/page/DS09ZUmdOeGVKUGpn}
    \item GitHub Issue：\url{https://github.com/econsun/XMU-AdvMacro-II/issues}
\end{itemize}

\vspace{\baselineskip}
\begin{flushright}
Rongqing SUN\\
福建\hspace{0.35em}\textperiodcentered\hspace{0.35em}厦门\\
2026 年 3 月 30 日
\end{flushright}

\clearpage
\pagestyle{symplain}



% 献上我的拙作
